% Standalone source for the Quanv1dLayer circuit figure.
% Edit this file, then recompile it to regenerate quanv_circuit.pdf:
%     pdflatex quanv_1D_layer.tex
% The main thesis includes the resulting PDF via \includegraphics.
\documentclass[border=8pt]{standalone}
\usepackage{tikz}
\usetikzlibrary{quantikz2}
\begin{document}
\begin{quantikz}[column sep=0.35cm, row sep=0.45cm]
    \lstick{$\ket{0}$} & \gate{R_P(\tilde{x}_1)}\gategroup[4,steps=4,style={dashed,rounded corners,inner xsep=4pt},background]{Layer $\ell$ (repeated $L=3$ times)} & \gate[wires=2]{ZZ} & \gate{\mathrm{Rot}(\theta^{(\ell)}_1)} & \ctrl{1} & \gate{\mathrm{Rot}(\phi^{\text{lens}}_1)}\gategroup[4,steps=1,style={dotted,rounded corners,inner xsep=4pt},background]{Quantum Lens} & \meter{Z, X, ZZ} \\
    \lstick{$\ket{0}$} & \gate{R_P(\tilde{x}_2)} & & \gate{\mathrm{Rot}(\theta^{(\ell)}_2)} & \targ{} & \gate{\mathrm{Rot}(\phi^{\text{lens}}_2)} & \meter{Z, X, ZZ} \\
    \lstick{$\ket{0}$} & \gate{R_P(\tilde{x}_3)} & \gate[wires=2]{ZZ} & \gate{\mathrm{Rot}(\theta^{(\ell)}_3)} & \ctrl{1} & \gate{\mathrm{Rot}(\phi^{\text{lens}}_3)} & \meter{Z, X, ZZ} \\
    \lstick{$\ket{0}$} & \gate{R_P(\tilde{x}_4)} & & \gate{\mathrm{Rot}(\theta^{(\ell)}_4)} & \targ{} & \gate{\mathrm{Rot}(\phi^{\text{lens}}_4)} & \meter{Z, X, ZZ}
\end{quantikz}
\end{document}
% Note: R_P denotes a dynamically selected Pauli rotation (R_Y, R_Z, or R_X) depending on layer depth mod 3. The ring connections for ZZ and CNOT are simplified for visual clarity.